\apendice{Documentación de usuario}

\section{Introducción}

En este manual se incluye una guía destinada a que el usuario final conozca el funcionamiento de la \eng{web} PythIA. Concretamente se van a describir los requisitos que debe cumplir el usuario para utilizar la aplicación (ver sección~\ref{usu:req}) así como el método de instalación de la aplicación para el usuario final (ver sección~\ref{usu:inst}). Finalmente, se incluye un manual explicando todas las funcionalidades de la aplicación (ver sección~\ref{usu:manual}).

\section{Requisitos de usuarios}
\label{usu:req}

Para utilizar la aplicación PythIA, el usuario únicamente necesita disponer de acceso a Internet y de un navegador \eng{web} moderno compatible con los estándares actuales~\cite{usuario:webStandards}. Se recomienda utilizar un navegador con visor \file{PDF} integrado para visualizar o descargar los documentos \file{PDF} de las licitaciones, 

La aplicación ha sido probada satisfactoriamente en los navegadores \file{Google Chrome}, \file{Mozilla Firefox} y \file{Microsoft Edge}. Concretamente se han probado en las siguientes versiones de los navegadores:
\begin{itemize}
\tightlist
	\item \file{Google Chrome} 137.0 a  148.0.
	\item \file{Mozilla Firefox} 139.0 a 150.0.
	\item \file{Microsoft Edge} 137.0 a 148.0.
\end{itemize}

Para acceder a las funcionalidades del sistema es necesario tener una cuenta de usuario registrada y autenticarse mediante correo electrónico y contraseña.

Respecto al \eng{hardware} no existen requisitos específicos, puesto que PythIA ha sido desarrollada siguiendo un diseño \eng{responsive}. Esto permite su utilización desde diferentes tipos de dispositivos, incluyendo ordenadores de sobremesa, ordenadores portátiles, tabletas y teléfonos móviles, adaptando automáticamente la interfaz al tamaño de pantalla disponible. 

\section{Instalación}
\label{usu:inst}

No es necesario instalar ningún \eng{software} adicional, ya que la aplicación se ejecuta a través del navegador \eng{web}. Solo se tiene que acceder a la dirección \url{https://pythia.es/} e iniciar sesión con sus credenciales. En caso de no disponer de una cuenta, deberá registrarse.

Las actualizaciones de la aplicación se realizan de forma centralizada en el servidor, por lo que los usuarios siempre acceden a la versión más reciente sin necesidad de efectuar ninguna acción adicional en sus dispositivos.

\section{Manual del usuario}
\label{usu:manual}

En este manual se describen las funcionalidades que pueden llevarse a cabo en la aplicación. Se va a explicar tanto los elementos generales del sistema como cada una de las páginas de la aplicación describiendo sus componentes y funcionalidades. Además se va a explicar la navegabilidad.

\subsection{Página de inicio}
\label{manual:inicio}

Es la primera página que se ve al acceder a la página \eng{web} de \href{https://pythia.es/}{PythIA} (ver imagen~\ref{fig:Manual_usuario/Pagina_inicio})

\imagendiagram{Manual_usuario/Pagina_inicio}{Página de inicio de la aplicación \eng{web} PythIA.}{1.2}

Desde ella se puede acceder a la página de <<Inicio de sesión>> (ver sección~\ref{manual:inicioSesion}) y a la página de <<Registro>> (ver sección~\ref{manual:registro})

\subsection{Inicio de sesión}
\label{manual:inicioSesion}

Si el usuario ya tiene cuenta en la aplicación debe acceder a está página para introducir su correo electrónico y su contraseña y acceder a las funcionalidades de la aplicación. 

Desde esta página se puede acceder a la <<Página principal>> (ver sección~\ref{manual:pagPrincipal}), <<Registro>> (ver sección~\ref{manual:registro}), <<Recuperar contraseña>> (ver sección~\ref{manual:recuperar}) y <<Página de inicio>> (ver sección~\ref{manual:inicioSesion}).

Además permite ver un tutorial interactivo que muestra las funcionalidades de la página.

\subsection{Recuperar contraseña}
\label{manual:recuperar}

\subsection{Registro}
\label{manual:registro}

\subsection{Página principal}
\label{manual:pagPrincipal}

\subsection{Página de consultas}
\label{manual:pagConsulta}

\subsection{Comparar modelos}
\label{manual:compararModelos}

\subsection{Evaluar modelos}
\label{manual:evaluar}

\subsection{Historial de consultas}
\label{manual:historial}

\subsection{Perfil de usuario}
\label{manual:perfil}

\subsection{Estadísticas}
\label{manual:estad}

\subsection{Administrar documentos}
\label{manual:documentos}

Solo administradores

\subsection{Administrar usuarios}

Solo administradores

\subsection{Navegabilidad}